\documentclass[12pt]{article}
\usepackage{threeparttable}
\usepackage{booktabs}
\usepackage{inputenc, amsmath,amssymb, amsthm}
\usepackage[font=small,labelfont=bf,
   justification=justified,
   format=plain]{caption}
\usepackage{subcaption}
\usepackage{longtable}
\usepackage{adjustbox}
\usepackage{graphicx}
\newtheorem{theorem}{Theorem}
\usepackage{xcolor}
\usepackage{float}
\usepackage[colorlinks=true,
            linkcolor=blue,      % color of internal links (sections, pages, etc)
            citecolor=blue,      % color of citation links
            urlcolor=blue,       % color of URL links
            filecolor=blue]{hyperref}
\usepackage{url}
\usepackage{array}
\usepackage{url}
\usepackage{tikz}
\usetikzlibrary{shapes.geometric, arrows, positioning, fit}
\usetikzlibrary{calc}
\usepackage{titling}
\usepackage{lipsum}
\usepackage{natbib}
\usepackage{setspace}
\usepackage[margin=.75in]{geometry}
\usepackage{xcolor}
\usepackage{amsmath, amssymb, amsfonts}  % loads all AMS symbols & environments

\theoremstyle{definition}
\newtheorem{definition}{Definition}[section]
\usepackage[english]{babel}

\onehalfspacing


\graphicspath{{../../../../../Dropbox/Lawsuits Asymmetry/output}}


\title{Streamlining vs Density Bonuses: \\
 What Drives Affordable Housing Development?}
\author{Tyler Jacobson}
\begin{document}
\maketitle
\section{ED1 Overview and Modifications}
Karen Bass issued the original \hyperlink{https://mayor.lacity.gov/news/mayor-bass-signs-executive-directive-dramatically-accelerate-and-lower-cost-affordable-housing}{ED1} in December of 2016. The program was intented to dramatically speed up the approval process for 100 percent affordable housing and to make such approvals ministerial instead of discretionary. The key provisions:
\begin{enumerate}
    \item ``Applications for 100\% affordable housing projects[\dots] shall be, and hereby are deemed exempt from discretionary review processes [\dots] as long as such plans do not require any zoning change, variance, or General Plan amendment.'' 
    \item ``An application for the development of a 100 percent affordable housing project or Shelter may use the density permitted for that site either by the applicable zoning or the General Plan Land Use Designation, consistent with state law. In addition, a project may utilize the State Density Bonus and LAMC bonuses, incentives, waivers and concessions if such are in compliance with the applicable requirements.'' 
    \item ``I further direct all applicable City Departments to process clearances and utility releases related to building permit applications, certificates of occupancy, or temporary certificates of occupancy within 5 business days for 100 percent affordable housing projects and within 2 business days for Shelters.''
\end{enumerate}

\subsection{Changes to ED1}
ED1 has gone through a variety of changes due to pushback. In June 2023, Bass exempts single family areas. The most recent \hyperlink{https://planning.lacity.gov/odocument/4bdff0d5-a458-4bcc-a8c5-451e4af45ea7/ED1_revised_memo_3.pdf}{version} from July 2024 (prior to the permanent version the policy) adds the following provisions:
\begin{enumerate}
    \item ``and in no instance shall the project be located in a single family or more restrictive zone''
    \item ``The project site does not include any parcels located in a manufacturing zone that does not allow multifamily residential uses.''
    \item ``The project does not include any parcels that are included in the National Register of Historic Places or the California Register of Historical Resources, either individually or within a historic district, or included within a Historic Preservation Overlay Zone (HPOZ), or designated as a City Historic-Cultural Monument''
    \item ``Projects seeking Density Bonuses under LAMC Section 12.22A.25 shall be eligible for no more than five incentives and one waiver.'' 
    \item ``A project in a residential land use designation shall be eligible to request no more than a 100 percent increase in floor area, or up to a floor area ratio of 3.5 to 1, whichever is greater.''
    \item ``A project in a residential zone shall be eligible to receive no more than a total project height increase of three stories, or 33 feet, in excess of the otherwise applicable height limit imposed by the project's zoning.''
    \item ``The project is not located on a parcel or parcels subject to the Rent Stabilization Ordinance (RSO) containing 12 or more total units that are occupied or were occupied in the five-year period preceding the application.''
\end{enumerate}

Commentary on the modification from a KCRW \hyperlink{https://www.kcrw.com/shows/kcrw-features/stories/ed1-eviction-affordable-housing}{article}:
\begin{itemize}
    \item ``Many developments have been deemed “out of scale” by neighborhood groups and LA City planners. Projects currently in the ED1 pipeline would demolish at least several hundred rent-stabilized units, displacing low-income renters like Juarez.''
    \item Last summer, under pressure from homeowner groups, Mayor Bass modified ED1 to declare single-family home zones ineligible for fast-tracked low-income housing. That change exempted more than 70\% of LA from the affordable housing fast track. Bass’ newest revisions also exempt local historic districts.
\end{itemize}

\subsection{Permanent ED1}
In December of 2025, the LA city council voted to make the ED1 program permanent as the LA Times \hyperlink{https://www.latimes.com/california/story/2025-12-09/la-city-council-passes-ordinance-to-streamline-affordable-housing}{reports}. The permanent policy appears to mantain the reforms that Bass implemented in the 2024 version of the policy.


\end{document}
